% =====================================================================
%  The Radiomics--Microbiome Axis Layer
%  Manuscript-style paper (two-column article, zero external .cls dep).
%  Restructured 2026-05-18 from the single-column proposal report,
%  which is preserved unmodified at:
%     docs/proposal/archive/report_sanomap_radiomics_layer.tex
%  Build:  pdflatex paper_sanomap_radiomics_layer.tex   (run twice)
% =====================================================================
\documentclass[10pt,twocolumn,letterpaper]{article}

\usepackage[utf8]{inputenc}
\usepackage[T1]{fontenc}
\usepackage{geometry}
\geometry{margin=0.75in, columnsep=0.28in}
\usepackage{titlesec}
\usepackage{enumitem}
\usepackage{booktabs}
\usepackage{amsmath}
\usepackage{amssymb}
\usepackage[hidelinks]{hyperref}
% microtype intentionally omitted: basictex Type1 fonts are non-scalable,
% so pdftex font expansion is fatal. Reliability > typographic polish here.

% Compact list spacing for a dense two-column manuscript
\setlist{nosep,leftmargin=1.2em}

% Journal-style section headings
\titleformat{\section}
  {\normalfont\large\bfseries}{\thesection.}{0.5em}{}
\titleformat{\subsection}
  {\normalfont\normalsize\bfseries}{\thesubsection}{0.5em}{}
\titlespacing*{\section}{0pt}{1.4ex plus 0.4ex}{0.8ex}
\titlespacing*{\subsection}{0pt}{1.1ex plus 0.3ex}{0.6ex}

\setlength{\parindent}{1em}
\setlength{\parskip}{0.15ex}

\title{\vspace{-3ex}\textbf{The Radiomics--Microbiome Axis Layer:\\
A Hybrid Text--Vision Pipeline with Five-Gate\\
Acceptance for Knowledge-Graph Extension}}

\author{Junseong Lee\\
\small SanoMap Project --- Radiomics-Microbiome Axis Layer\\
\small \texttt{junseong.lee652@gmail.com}}

\date{\small Draft of \today\ \textbf{(pre-print; benchmark metrics gated, see \S\ref{sec:eval})}}

\begin{document}

\twocolumn[
\begin{@twocolumnfalse}
\maketitle
\begin{abstract}
\noindent
\textbf{Background.} Literature-mined microbiome knowledge graphs such as
the SanoMap MINERVA platform connect microbes directly to diseases but omit
the quantitative, tissue-level imaging phenotypes that mechanistically link
them. Much of the quantitative microbe--imaging correlation evidence is
additionally locked inside figures, invisible to text-only pipelines.
\textbf{Methods.} We build a radiomics-first intermediary layer that inserts
\texttt{RadiomicFeature} and \texttt{BodyCompositionFeature} nodes between
microbe and disease. A two-track pipeline (a text track over 1{,}016
open-access papers; a vision track that reads correlation heatmaps with a
vision-language model) feeds a five-gate acceptance model: dense-retrieval
candidate generation, UMLS type-grounded entity sanitization,
self-consistency relation acceptance, dual-verifier vision consensus behind a
deterministic pre-verifier gate chain, and a stratified gold-label benchmark.
Accepted edges are reconciled into a single provenance-stamped graph export
and loaded into Neo4j with a read-only query layer.
\textbf{Results.} The graph carries 74 phenotype--disease
\texttt{ASSOCIATED\_WITH} edges, 29 signed microbe--disease edges, and---after
a retroactive vision audit and a graph reconciliation that composed both the
vision retraction and the UMLS drop onto one base---7
\texttt{CORRELATES\_WITH} microbe--feature edges (1 vision, 6 text), closing
62 traversable microbe$\rightarrow$feature$\rightarrow$disease paths a
single-hop schema cannot express. The vision track's measured contribution is
deliberately scoped to one audited case-study figure plus a
hallucination-resistant gating methodology. Precision/recall and
intra-annotator Cohen's $\kappa$ are gated on a non-negotiable 14-day
temporal re-labeling window and are reported as pending, not estimated.
\textbf{Conclusion.} Making the imaging phenotype explicit is feasible,
reviewable, and reproducible; the engineering contribution is the
acceptance-gate architecture and an honest accounting of where each evidence
modality does and does not deliver.

\vspace{1ex}
\noindent\textbf{Keywords:} knowledge graph; radiomics; gut microbiome;
vision-language models; biomedical relation extraction; Neo4j; evidence
provenance.
\end{abstract}
\vspace{3ex}
\end{@twocolumnfalse}
]

% =====================================================================
\section{Introduction}
% =====================================================================

Biomedical knowledge graphs derive their clinical value from the
completeness of the evidence they encode. Transformer-based language models
and generative AI have enabled large-scale construction of such graphs by
systematically mining verified literature. The SanoMap project's MINERVA
platform (Microbiomes Network for Research and Visualization Atlas)
demonstrates this well: it maps associations between gut microbes and
systemic diseases at scale, producing a clinically navigable graph from
thousands of curated papers.

Yet MINERVA's graph encodes a structurally incomplete picture. It connects
microbes directly to diseases but omits the quantitative, tissue-level
intermediate phenotypes that explain the biological mechanism linking them.
Radiomics and body-composition imaging provide exactly these intermediaries.
Skeletal muscle index, visceral adipose tissue fraction, liver fat fraction,
and GLCM texture signatures are extracted directly from CT and MRI and are
established biomarkers for sarcopenia, nonalcoholic fatty liver disease, and
colorectal cancer. The literature increasingly demonstrates that gut
microbiome composition measurably influences these specific tissue-level
phenotypes---making them natural intermediary nodes in any complete
microbiome--disease graph.

A second structural gap concerns the evidence modality itself. Text-mining
pipelines extract findings authors highlight in prose, but much of the
quantitative correlation data linking microbial taxa to imaging feature
values is encoded in figures. Gradient correlation heatmaps display matrices
of $r$-values linking dozens of taxa to radiomic measurements---data NLP
alone cannot parse. This is a systematic evidence loss in text-only
pipelines.

This work addresses both gaps. The central research questions are:
\textbf{(RQ1)} Can a hybrid pipeline combining large language models and
vision-language models, with deterministic verification, reliably extract
quantitative correlation coefficients directly from medical figures? And
\textbf{(RQ2)} does adding radiomics as an ontological layer---filling the
microbe$\rightarrow$imaging-phenotype$\rightarrow$disease path---meaningfully
increase the clinical completeness of the MINERVA knowledge graph? A third,
engineering question runs underneath both: \textbf{(RQ3)} what acceptance
architecture keeps such a hybrid graph \emph{reviewable}---auditable
edge-by-edge---rather than merely large?

This paper's contribution is not a claim of upstream parity. It is (i) an
explicit intermediary-layer ontology, (ii) a five-gate acceptance model that
converts a proof-of-concept pipeline into a measurable artifact, (iii) a
retroactive audit that honestly rescopes the vision track after a structural
verifier weakness was found, and (iv) a single reconciled, provenance-stamped
graph reproducibly loadable into Neo4j.

% =====================================================================
\section{The Intermediary Layer: Concept and Schema}
% =====================================================================

\subsection{Core idea}

MINERVA's current graph encodes
\texttt{(Microbe)$\rightarrow$[ASSOCIATED\_WITH]$\rightarrow$(Disease)}. This
edge is clinically underspecified. The biologically accurate path is

\begin{center}
\small
\texttt{Microbe} $\rightarrow$ \textit{correlates with} $\rightarrow$
\texttt{Radiomic/BodyComp\-Feature} $\rightarrow$ \textit{associates with}
$\rightarrow$ \texttt{Disease}.
\end{center}

This three-node path is documented concretely: \textit{Prevotella
nigrescens} correlates with GLCM Correlation texture on CT at $r=0.95$
(PMC10605408); cirrhosis dysbiosis enriches Proteobacteria alongside
measurable changes in liver surface nodularity and visceral fat;
\textit{Bifidobacterium} depletion co-occurs with elevated visceral adipose
tissue and downstream obesity. These are quantitatively measured, directly
evidenced relations that a text-only microbe$\rightarrow$disease edge
collapses into a single mechanism-blind link. Making the intermediary
explicit enables queries impossible in MINERVA's schema, e.g.\ ``what CT
texture signature mediates the association between gut dysbiosis and
sarcopenia?''

\subsection{Node and edge ontology}

Node types extend MINERVA: \texttt{RadiomicFeature} (texture, morphological,
intensity features---GLCM, wavelet, first-order, shape);
\texttt{BodyCompositionFeature} (sarcopenia, visceral adipose tissue,
skeletal muscle index, liver fat fraction, myosteatosis); \texttt{BodyLocation};
\texttt{ImagingModality}; \texttt{ImageRef} (verified figure references with
Vision Track provenance); \texttt{Microbe} / \texttt{MicrobialSignature}; and
\texttt{Disease}.

The direct-evidence edge policy is deliberately evidence-typed:
quantitatively verified microbe$\rightarrow$feature edges (a vision-confirmed
$r$-value) use \texttt{CORRELATES\_WITH}; text co-mention
feature$\rightarrow$disease edges use \texttt{ASSOCIATED\_WITH}; signed
microbe$\rightarrow$disease edges use
\texttt{POSITIVELY\_/NEGATIVELY\_CORRELATED\_WITH}; structural edges use
\texttt{MEASURED\_AT}, \texttt{ACQUIRED\_VIA}, \texttt{REPRESENTED\_BY}. The
\texttt{CORRELATES\_WITH}/\texttt{ASSOCIATED\_WITH} distinction preserves
evidence-level semantics across the graph and is never collapsed. Within-paper
bridge matches that only share disease context are written to an audit
artifact as hypotheses and are \emph{never} ingested as asserted edges.

% =====================================================================
\section{Methods}
% =====================================================================

\subsection{Data sourcing}

Data was sourced from PubMed Central Open Access via a multi-lane query
strategy over six clinical domains where the microbiome--radiomics link is
best evidenced: (1) strict radiomics + microbiome; (2) body composition +
microbiome; (3) liver radiomics + microbiome; (4) bone-density (DXA) +
microbiome; (5) lung-CT phenotypes + microbiome; (6) colorectal imaging +
microbiome. The corpus spans \textbf{1{,}016 papers} (640 from the initial
three lanes plus 376 net-new from lanes 4--6), with open-access full text
retrieved where a PMCID exists. Review, meta-analysis, and protocol papers
are excluded as non-direct-evidence targets.

\subsection{Two-track pipeline}

The \textbf{Text Track} extracts phenotype--disease associations at scale. A
vocabulary-matched stage maps imaging-phenotype terminology to canonical
names; microbial NER uses \texttt{d4data/biomedical-ner-all} on Apple Silicon
MPS, disease NER uses the \texttt{en\_ner\_bc5cdr\_md} spaCy pipeline; a
shared span-cleanup layer rejects wordpiece artifacts, generic placeholders,
and clause-prefix fragments before relation classification.

The \textbf{Vision Track} addresses the figure-data gap. A vision-language
model proposes $r$-values for microbe$\rightarrow$feature correlations read
from heatmap figures; each proposal is deterministically cross-checked
against the figure's gradient color-scale legend at pixel level.

\subsection{Five-gate acceptance model}
\label{sec:gates}

Tracks alone do not establish research-defensibility; structurally distinct
gates do. A candidate edge enters the graph only when every gate that applies
to its evidence track passes.

\begin{enumerate}
\item \textbf{Retrieval (text).} A dense-retrieval index over
\texttt{BioClinical-ModernBERT-base} embeddings (mean-pooled, L2-normalized;
FAISS where available, NumPy cosine fallback at $\sim$30k sentences) replaces
the original hand-curated 25-alias substring filter, removing its recall
ceiling on paraphrastic mentions (``loss of muscle mass at L3''). Per-feature
similarity threshold $\tau$ is calibrated per concept against labeled data.

\item \textbf{Entity sanitization (UMLS TUI grounding).} Microbial surface
forms must ground to a microbe-class UMLS Semantic Type at similarity
$\geq 0.85$; accept set $\{$T005 Virus, T007 Bacterium, T194 Archaeon, T204
Eukaryote$\}$ with an explicit non-microbe model-organism deny-list. A live
audit of the 8 accepted text-derived microbe$\rightarrow$feature edges drops
25\% (2/8): the entity-type error \texttt{gut bacterial ClpB-like gene
function} (similarity 0.799 $<$ 0.85, grounded to ``Bacteria'' T007) and the
typo \texttt{bacteriodetes} (no UMLS match). The 25\% drop sits below the
30\% recalibration threshold; the 0.85 floor is retained.

\item \textbf{Relation acceptance (text).} Gemini 2.5 Flash-Lite,
7-sample temperature-varied self-consistency, full agreement required;
measured full-consistency rate on the expanded corpus is 0.916.

\item \textbf{Verification (vision) --- dual verifier behind a
pre-verifier gate chain.} See \S\ref{sec:visionaudit}; this gate's design is
itself a contribution.

\item \textbf{Evaluation (gold benchmark).} A 5-stratum sampler
(\texttt{accepted\_edge}, \texttt{gemini\_rejected}, \texttt{vocab\_excluded},
\texttt{recall\_probe}, \texttt{random\_co\_occurrence}) drawn deterministically
from the production candidate space yields a 66-row gold set. The schema
defines a 6-class primary label, three secondary labels, eight numbered
edge-case decisions, and a 14-day intra-annotator temporal-relabeling
protocol with Cohen's $\kappa \geq 0.80$ (binary collapse) acceptance.
Pass-1 used computer-aided manual annotation (LLM proposer, single human
reviewer who overrode 7 rows under an imaging-derived scope rule); Pass-2 is
the annotator's independent re-label after the temporal window. Metrics that
depend on Pass-1+Pass-2 are pending (\S\ref{sec:eval}).
\end{enumerate}

\subsection{Graph integration and reproducibility}
\label{sec:neo4j}

Earlier iterations emitted several divergent CSV vintages and a hand-runnable
Cypher script, with no single coherent ``current graph.'' This release adds a
deterministic reconciliation step: one command reads the authoritative
assembled artifacts, applies the documented post-audit drop-list (the
retracted Firmicutes vision edge; the two UMLS-failed
microbe$\rightarrow$feature edges), and emits a single
\texttt{graph\_export/} bundle---\texttt{nodes.csv},
\texttt{relationships.csv}, an idempotent \texttt{import.cypher}, and a
\texttt{manifest.json} recording source-artifact vintages, the git commit,
pre/post counts, and every drop with its reason. This step immediately earned
its keep: it found that the pre-existing prose headline (``8
\texttt{CORRELATES\_WITH}: 1 vision + 7 text'') had applied only the vision
retraction to the 9-edge superset and never composed in the separately
recorded UMLS drop; composing both onto one base yields the coherent 7 (1
vision + 6 text). The bundle loads into Neo4j through a thin driver loader; a
read-only query module exposes the canonical traversals (three-hop microbe$\rightarrow$feature$\rightarrow$disease;
feature-by-disease; signed microbe--disease; feature-by-location; full
modality-backbone chain; vision-verified edges). Provenance, not just the
graph, is the deliverable: every edge resolves back to a PMID, sentence, or
figure panel, and every count in this paper is reproducible from the
manifest.

% =====================================================================
\section{Results}
\label{sec:results}
% =====================================================================

\subsection{Extracted graph}

Across the 1{,}016-paper corpus the pipeline produced:

\begin{itemize}
\item \textbf{5{,}721} phenotype text mentions across 30 clinically
meaningful disease targets (initial 640-paper corpus; expanded-corpus
re-extraction in progress);
\item \textbf{74} \texttt{ASSOCIATED\_WITH} phenotype$\rightarrow$disease
edges (Gemini full-consistency 0.916 on the expanded corpus);
\item \textbf{29} signed microbe--disease edges (14 positive, 15 negative)
spanning cirrhosis, NAFLD, colorectal cancer, IBD, celiac disease, chronic
HCV, and SSc-associated pulmonary fibrosis;
\item \textbf{7} \texttt{CORRELATES\_WITH} microbe$\rightarrow$feature edges
\emph{after} reconciliation (\S\ref{sec:visionaudit}): 1 vision
(\textit{Prevotella nigrescens} $\rightarrow$ GLCM\_Correlation, $r=0.95$,
PMC10605408\_g004; pixel distance 0.05, support fraction 1.0) and 6 text
(Gemini 7/7);
\item \textbf{18} \texttt{BodyLocation} and \textbf{5}
\texttt{ImagingModality} nodes;
\item \textbf{189} reconciled relationship rows over \textbf{99} nodes. The
pre-reconciliation export reported 191 rows / 9 \texttt{CORRELATES\_WITH}; the
reconciliation step (\S\ref{sec:neo4j}) drops exactly two edges---one vision
retraction, one UMLS-failed entity---each recorded in the manifest with its
cited source artifact, rather than silently changing the headline.
\end{itemize}

Three surviving \texttt{CORRELATES\_WITH} edges close end-to-end three-hop
paths through features that also carry downstream \texttt{ASSOCIATED\_WITH}
disease edges---\textit{Ruminococcus}$\rightarrow$sarcopenia (37 diseases),
\textit{Peptostreptococcus stomatis}$\rightarrow$skeletal\_muscle\_index (7),
\textit{Eubacterium}$\rightarrow$visceral\_adipose\_tissue (18)---yielding
\textbf{62 traversable microbe$\rightarrow$feature$\rightarrow$disease
paths}, the query pattern MINERVA's single-hop schema cannot express. This is
direct evidence for RQ2: the intermediary layer adds traversable mechanism,
not just nodes.

\subsection{Corpus-level findings}

Two findings merit note. \textit{Faecalibacterium} depletion emerged as a
convergent signal across six distinct disease associations in the new liver
and colorectal lanes (chronic liver disease, colorectal cancer, IBD, celiac
disease, IBS, SSc-associated pulmonary fibrosis), consistent with its
established role as a butyrate-producing protective commensal.
\textit{Fusobacterium nucleatum} enrichment in colorectal cancer was
independently confirmed across four PMIDs from the colorectal lane,
validating that lane's signal quality.

\subsection{Vision-track audit and honest scoping}
\label{sec:visionaudit}

A retroactive audit of all vision inputs (13 current proposals + 1 historical
graph edge) revealed that dual-verifier consensus \emph{alone} is
structurally insufficient: the pixel HSV verifier and the independent VLM
verifier both consume the proposer's bounding box and proposed $r$-value, so
a self-consistent fabrication (proposer hallucinates both fields so the
predicted color matches the colorbar at that location on the assumed default
$\pm1.0$ scale) passes AND-consensus silently. The pixel verifier also
defaulted its observed range to $[-1,+1]$ rather than reading tick labels, so
log-fold-change ($\pm1.5$) heatmaps were validated as if Pearson $r$.

Three deterministic pre-verifier gates were added in
\texttt{src/vision\_gates.py} and run \emph{before} consensus: a
\textbf{caption gate} (requires Pearson/Spearman vocabulary; auto-fails
\texttt{log fold change}, \texttt{LFC}, \texttt{z-score}, \texttt{differential
abundance}); a \textbf{colorbar-detection gate} (a gradient legend must be
present); and a \textbf{range-sanity gate}
($|r_{\text{prop}}|\leq\max(|c_{\min}|,|c_{\max}|)+0.05$, with
$(c_{\min},c_{\max})$ read from the figure's tick labels by a focused VLM
call). On the audit batch ($n{=}14$) the chain rejected 6 figures (3 caption,
2 range-sanity, 1 colorbar); the dual verifier flagged 3 of the remaining 8
for human review. The historical \textit{Prevotella nigrescens} edge was
retained (caption-confirmed Spearman heatmap, extracted $-1$ to $+1$
colorbar, pixel distance 0.05, direct image inspection). The historical
\textit{Firmicutes}$\rightarrow$Total fat \% edge ($r=-0.95$, PMC6178902) was
\textbf{retracted}: the asterisked cell is in the figure's positive (red)
hemisphere, the colorbar runs $-1.5$ to $+1.5$ (a log-fold-change scale, not
Pearson $r$), and the prior ``pixel-verified'' pass was a circular
consistency check on a hallucinated bounding box.

This reframes the vision contribution honestly. Rather than a pre-audit ``9
edges (2 vision + 7 text)'' framing, the reconciled graph carries 7 (1 vision
+ 6 text), once the vision retraction and the UMLS drop are composed onto one
base (\S\ref{sec:neo4j}). The publishable vision result is therefore
\emph{not} edge volume;
it is (a) one fully audited case-study figure and (b) a hallucination-resistant
gating methodology with a documented residual failure mode---wrong-sign
hallucination when the bounding box happens to point at a same-color cell
elsewhere---whose cheapest mitigation (a sign-check gate comparing claimed
sign vs.\ VLM-observed color hemisphere) is specified for future work. This is
the honest answer to RQ1: deterministic verification is necessary but not
sufficient; the gate chain, not the consensus, is what makes vision evidence
defensible.

\subsection{Verification and software quality}

The pipeline carries 280 automated test checks across the full suite (up from
156 before the structural-upgrade work; +105 new tests cover the UMLS
validator, dense retrieval, dual verifier, gold-set sampler, and benchmark
evaluator), with the graph-export reconciliation and query layer added in
this release additionally covered. No regressions on prior tests.

\subsection{Gated metrics}
\label{sec:eval}

Measured precision/recall and intra-annotator Cohen's $\kappa$ are
deliberately reported as \textbf{pending}, not estimated. The benchmark
requires a non-negotiable 14-day intra-annotator temporal window between
Pass-1 (closed 2026-05-07) and Pass-2 (earliest 2026-05-21); a shorter
interval leaks short-term memory and inflates $\kappa$ artificially.
Reporting an estimate here would violate the evidence-typed discipline the
rest of the system enforces. On the window's close we will report binary
precision/recall/F1 with 95\% Wilson confidence intervals (wider, $\sim\pm
0.10$, because the corpus-limited gold set is 66 rows, not the 150-row
design) and Cohen's $\kappa$ on the two human passes.

% =====================================================================
\section{Discussion: Comparison with MINERVA}
% =====================================================================

The system extends MINERVA, not replicates it. Each gap below is an
explicitly acknowledged design decision, not an error.

\noindent\textbf{Gap 1 --- NER models.} MINERVA fine-tuned a custom
DistilBERT (BNER2.0) on a proprietary corpus (F1$=$0.914; weights
unreleased). We use \texttt{d4data/biomedical-ner-all} (microbe) and
\texttt{en\_ner\_bc5cdr\_md} (disease), selected via a controlled 5-paper
benchmark. Residual risk: lower precision on phylum/genus/species
disambiguation, partially compensated by span cleanup and a genus allowlist.
Closure paths: Professor-mediated BNER2.0 access, or the public
\texttt{OpenMed-NER-SpeciesDetect} (DeBERTa-v3+LoRA, SotA on Linnaeus/S-800)
as a no-fine-tune upgrade.

\noindent\textbf{Gap 2 --- Relation model.} MINERVA fine-tuned
BioMistral-AUG-7B (weights unreleased). We use Gemini 2.5 Flash-Lite
zero-shot with 7-sample self-consistency (full-consistency rate 0.896--0.916,
a calibrated reliability measure, not a P/R claim). Closure path: fine-tune
BiomedBERT on MicrobioRel (22 gut-microbiome relation types) or the Lever
et al.\ 2023 labeled set.

\noindent\textbf{Gap 3 --- UMLS normalization: closed and exceeded.} We match
MINERVA's CUI-linked deduplication (ScispaCy \texttt{en\_core\_sci\_lg},
2022 UMLS KB; 96\% / 83\% CUI coverage on the initial / expanded corpora) and
go beyond it by promoting the same CUI/TUI metadata to an \emph{active
acceptance gate}: a microbe failing to ground to a microbe-class TUI at
similarity $\geq 0.85$ is rejected pre-classification with an audited
\texttt{drop\_reason}.

\begin{table}[t]
\centering
\small
\caption{Gap status relative to MINERVA.}
\label{tab:gap}
\begin{tabular}{@{}p{0.30\columnwidth}p{0.13\columnwidth}p{0.42\columnwidth}@{}}
\toprule
\textbf{Gap} & \textbf{Sev.} & \textbf{Status} \\
\midrule
NER (microbe) & Med & Substitute identified; TUI gate mitigates type errors \\
NER (disease) & Low & BC5CDR competitive; SotA upgrade identified \\
Relation model & Med & Self-consistency 0.916; fine-tune path documented \\
UMLS norm. & High & \textbf{Closed \& exceeded} (TUI-gated, audit-active) \\
Text recall ceiling & Med & \textbf{Closed} (ModernBERT dense retrieval) \\
Vision monoculture & High & \textbf{Partially closed}: dual gate alone
insufficient; 3 pre-verifier gates added; audit + reconciliation
9$\rightarrow$7 edges \\
Measured P/R & High & \textbf{Pending} Pass-2 (\S\ref{sec:eval}) \\
\bottomrule
\end{tabular}
\end{table}

% =====================================================================
\section{Limitations and Future Work}
% =====================================================================

\noindent\textbf{Radiomic parameter abstraction.} The pipeline captures
modality and body location but abstracts away bin width, voxel resampling,
slice spacing, and ROI method---rarely reported in mined text. Two GLCM
entropy values from different bin widths become one node. Future work: map to
IBSI feature-instance codes encoding extraction context.

\noindent\textbf{Demographic heterogeneity.} Body-composition metrics depend
strongly on age, sex, and BMI; the graph aggregates across heterogeneous
cohorts without a \texttt{PopulationContext} property. Future work: a cohort
node or edge stratum enabling demographic filtering.

\noindent\textbf{Inter-scanner variability --- a vision-track advantage.} By
extracting $r$-values from heatmaps rather than absolute radiomic values, the
vision track inherently normalizes across scanner hardware: a rank-order
correlation is far more stable across sites than an absolute HU value, giving
vision edges higher cross-site generalizability than value-anchored text
associations.

\noindent\textbf{Observational only.} The graph encodes association, not
causality. Future work: \texttt{Intervention} nodes (FMT, antibiotics,
probiotics, diet) with typed \texttt{MODIFIES}/\texttt{REVERSES} edges,
shifting toward a partially causal model---requiring longitudinal
pre/post datasets targeted in future harvest-lane design.

\noindent\textbf{Sparsity and track complementarity.} The three evidence
layers are complementary by construction. The text track
(phenotype$\rightarrow$disease) is a broad associative network (74 edges, 30
diseases); the text track (microbe$\rightarrow$feature) contributes 6
self-consistency-validated \texttt{CORRELATES\_WITH} edges; the vision track
contributes 1 quantitatively verified edge plus the gating methodology. The
microbe$\rightarrow$feature edges are what close the 62 three-hop paths the
intermediary-layer thesis requires; remaining \texttt{CORRELATES\_WITH} edges
without downstream disease edges will close as text-track phenotype
extraction expands to additional feature vocabularies. The graph is
intentionally sparse and reviewable rather than dense and unaudited---this is
a design stance, answering RQ3.

% =====================================================================
\section{Conclusion}
% =====================================================================

Inserting an explicit radiomics--body-composition intermediary between
microbe and disease is feasible, mechanistically meaningful, and---given the
five-gate acceptance model and a single provenance-stamped graph
export---reproducible and reviewable. The vision track's honest contribution
is a hallucination-resistant verification methodology and one audited
case study, not edge volume; the benchmark's headline metrics are gated on a
disciplined temporal window and reported as pending rather than estimated.
The system's defensibility comes from where it refuses to overclaim as much
as from what it extracts.

\vspace{1ex}
\noindent\textbf{Reproducibility.} Code, tests, the reconciled
\texttt{graph\_export/} bundle with its provenance manifest, the Neo4j
loader, and the read-only query layer are tracked in the project repository.
The pre-reframing single-column proposal report is preserved unmodified at
\texttt{docs/proposal/archive/}. Every numeric claim here is reproducible
from the export manifest; gated metrics will be appended on the temporal
window's close.

\end{document}
